\documentclass[11pt]{article}
\usepackage[margin=1in]{geometry}
\usepackage{booktabs}
\usepackage{amsmath,amssymb}
\usepackage{graphicx}
\usepackage{multirow}
\usepackage{xcolor}
\usepackage{hyperref}
\usepackage{natbib}
\usepackage{authblk}
\usepackage{subcaption}
\usepackage{tablefootnote}
\usepackage{enumitem}

\hypersetup{colorlinks=true, linkcolor=blue, citecolor=blue, urlcolor=blue}

\newcommand{\kstar}{K^*}
\newcommand{\thetaval}{\theta}

\title{Bounded Generational Discrimination in Recursive Synthetic Text:\\K* Saturates Within Three Passes}

\author{Anonymous}
\date{}

\begin{document}
\maketitle

\begin{abstract}
Web-scale training corpora increasingly contain synthetic text generated by large language models (LLMs), and this text can itself seed further rounds of generation, creating recursively contaminated data of unknown generational depth.
Current detection methods treat synthetic text as a binary category (human vs.\ machine), ignoring the progressive quality degradation across recursive generations.
We introduce \emph{generational discrimination depth} $\kstar$, defined as the maximum generation depth at which pairwise statistical discrimination remains reliable above a calibrated threshold $\thetaval$.
We provide an information-theoretic argument based on the data processing inequality (DPI) that $\kstar$ must be finite for any fixed generator, then conduct the first systematic empirical measurement of $\kstar$ across 9 model configurations (5 architecture families, 1B--14B parameters), 4 source domains (C4, Wikipedia, arXiv, Code/Python), and 4 decoding strategies.
Our central empirical finding is that $\kstar$ saturates rapidly: $\kstar \in \{1,2,3\}$ across all tested conditions.
Wikipedia yields $\kstar{=}1$ (6/6 models under scorer control), C4 web text yields $\kstar{=}2$ as the dominant mode (8/9 configurations), and structured academic text (arXiv) extends to $\kstar{=}3$ (5/5 domain-competent models across 4 architecture families).
Code/Python shows $\kstar{=}2$ for 2 of 5 tested models, with the remaining 3 at $\kstar{=}1$, indicating model-dependent behavior in intermediate-complexity domains.
Cross-model transfer retains 80--92\% accuracy, and scorer-invariant feature ordering ($\rho{>}0.85$) confirms depth signals are text-intrinsic.
Downstream filtering experiments show that binary removal of detected synthetic text provides a small but statistically significant improvement (+0.76pp MMLU), while graduated depth-aware weighting offers no additional benefit over binary filtering at any contamination level tested.
We discuss the implications of rapid $\kstar$ saturation for data governance and the practical limits of depth-aware curation.
\end{abstract}


%==============================================================================
\section{Introduction}
%==============================================================================

Training data for large language models now contains a growing fraction of machine-generated text.
As LLMs are deployed at scale for content creation, web corpora accumulate synthetic text that can feed into subsequent training rounds, producing recursively contaminated data of increasing generational depth.
This recursive contamination poses a qualitatively different challenge from the single-generation detection problem addressed by existing binary detectors.

The practical consequences of recursive contamination are becoming clear.
Model collapse studies \citep{shumailov2024ai,alemohammad2024selfconsuming,dohmatob2024tale} demonstrate that training on self-generated data leads to progressive quality degradation, distribution narrowing, and loss of tail knowledge.
Yet the detection literature treats all synthetic text identically: a passage is either human-written or machine-generated, with no distinction between first-generation and fourth-generation synthetic content.
This binary framing misses critical information about the severity of contamination.

We formalize this gap through the concept of \emph{generational discrimination depth} $\kstar$, defined as the deepest generation at which a classifier can reliably distinguish adjacent depth levels.
The data processing inequality (DPI) from information theory guarantees that $\kstar$ must be finite for any fixed generation process, since each generation pass is a lossy channel that cannot increase mutual information with the original source.
However, DPI provides no upper bound on $\kstar$: in principle, $\kstar$ could be large if each pass introduces distinctive artifacts.
Our core empirical contribution is the first systematic measurement of $\kstar$ across diverse conditions.

The central finding is striking: \textbf{$\kstar$ saturates within at most three generation passes across all conditions tested.}
This means recursive generation erases distinguishable depth structure extremely quickly.
For web text (C4), a single generation pass beyond $\kstar{=}2$ produces text that is statistically indistinguishable from deeper generations, an effect we term ``one-wash laundering'' for domains where $\kstar{=}1$.

We further find that $\kstar$ depends primarily on the source domain, not on model identity, scale, or decoding strategy.
Wikipedia text saturates at $\kstar{=}1$ across all 6 models tested (after controlling for scorer confounds), C4 web text at $\kstar{=}2$ (8/9 configurations), and structured arXiv text extends to $\kstar{=}3$ (5/5 domain-competent models).
Code/Python presents a more complex picture: 2 of 5 models achieve $\kstar{=}2$ while 3 remain at $\kstar{=}1$, suggesting that intermediate-complexity domains exhibit model-dependent behavior rather than a clean domain ceiling.

For downstream applications, we evaluate whether $\kstar$-aware data curation provides practical benefit beyond binary filtering.
The results are sobering: binary removal of detected synthetic text yields small improvements, but graduated depth-aware weighting provides \emph{zero} additional benefit over binary filtering at all contamination levels and on both domains tested.
This negative result for depth-aware curation does not diminish the characterization contribution of $\kstar$, which provides the first empirical evidence that recursive information loss is both rapid and bounded.

Our contributions:
\begin{enumerate}[nosep]
    \item We formalize generational discrimination depth $\kstar$ with an information-theoretic grounding (DPI), and conduct the first systematic empirical measurement across 9 model configurations, 4 domains, and 4 decoding strategies, finding $\kstar \in \{1,2,3\}$ in all tested conditions.
    \item We demonstrate that domain structural complexity is the primary determinant of $\kstar$: Wikipedia $\kstar{=}1$, C4/Code $\kstar \in \{1,2\}$, arXiv $\kstar{=}3$, with model-dependent behavior concentrated in intermediate-complexity domains.
    \item We provide a comprehensive evaluation of downstream $\kstar$-aware filtering, finding that binary filtering captures the available benefit while graduated depth weighting adds no incremental value, an important negative result for the field.
\end{enumerate}


%==============================================================================
\section{Related Work}
%==============================================================================

\paragraph{Model Collapse and Recursive Generation.}
\citet{shumailov2024ai} demonstrated that iterative training on model-generated data leads to progressive distribution collapse, with loss of tail diversity across generations.
\citet{alemohammad2024selfconsuming} formalized self-consuming generative models and characterized conditions under which iterative retraining preserves or degrades distributional properties.
\citet{dohmatob2024tale} provided theoretical analysis showing that model collapse can occur even with careful data mixing.
These works characterize the \emph{macroscopic} phenomenon of quality degradation but do not provide tools for \emph{detecting} or \emph{measuring} the generational depth of individual text passages.
Our work bridges this gap by operationalizing depth discrimination and measuring where it saturates.

\paragraph{Machine-Generated Text Detection.}
Binary detectors such as DetectGPT \citep{mitchell2023detectgpt}, DivEye, Binoculars, and Min-K\%++ achieve high accuracy for the human-vs-machine distinction.
\citet{drayson2025preventing} applied binary detection to importance resampling for preventing model collapse in pretraining data curation, representing the closest prior work to our filtering experiments.
All these methods perform binary classification and do not estimate generational depth.
Our 15-dimensional feature approach extends beyond binary detection to ordinal depth estimation, though we find the practical advantage is limited.

\paragraph{Data Selection and Curation.}
Distribution-matching methods such as DSIR \citep{xie2023data}, D4, and DoReMi select pretraining data to match a target distribution but do not account for recursive contamination depth.
Our filtering experiments (Section~\ref{sec:filtering}) compare $\kstar$-aware strategies against DSIR and unfiltered baselines, finding that binary filtering captures most of the available benefit.

\paragraph{Information Theory and Iterated Channels.}
The data processing inequality (DPI) states that for a Markov chain $X \to Y \to Z$, mutual information satisfies $I(X;Z) \leq I(X;Y)$.
Applied to recursive generation (each pass as a lossy channel), DPI guarantees that distinguishability between source and generated text can only decrease with depth, implying finite $\kstar$.
Our work provides the first empirical measurement connecting this theoretical guarantee to concrete $\kstar$ values.


%==============================================================================
\section{Method}
\label{sec:method}
%==============================================================================

\subsection{Recursive Depth Chain Generation}

Given a source corpus $\mathcal{D}_0$ of human-written documents, we construct a \emph{depth chain} $\{\mathcal{D}_0, \mathcal{D}_1, \ldots, \mathcal{D}_K\}$ where each $\mathcal{D}_{k+1}$ is produced by applying a language model $M$ to documents in $\mathcal{D}_k$.
We study two generation modes:

\begin{itemize}[nosep]
    \item \textbf{Continuation mode} (canonical): the model continues from a prefix of the source document. This isolates text-intrinsic degradation and provides the basis for cross-model claims.
    \item \textbf{Rewrite mode}: the model paraphrases the full source document under instruction. This amplifies depth signals but introduces mode-specific style artifacts that reduce cross-model transferability (61--63\% vs.\ 80--92\% retention).
\end{itemize}

All cross-model conclusions are based on continuation mode. Rewrite mode serves as an ablation.

\subsection{Feature Extraction}

We extract a 15-dimensional feature vector from each text sample, comprising:
\begin{itemize}[nosep]
    \item \textbf{Surprisal features} (6D): mean, standard deviation, and second-order differences of token-level surprisal under a reference model.
    \item \textbf{Entropy features} (4D): vocabulary entropy, bigram entropy, and their depth-differences.
    \item \textbf{Distributional features} (5D): type-token ratio, repetition rate, Zipf coefficient, and related statistics.
\end{itemize}

These model-free features were selected for cross-model universality. Cross-model feature ordering shows Spearman $\rho > 0.85$ (Section~\ref{sec:features}), confirming consistent discriminative structure across generator architectures.

\subsection{$\kstar$ Determination}
\label{sec:kstar-def}

For each adjacent depth pair $(k, k{+}1)$, we train a binary classifier (LightGBM) on the 15D features using 5-fold cross-validation and compute the pairwise AUC. We define:

\begin{equation}
    \kstar = \max\{k : \text{AUC}(k \text{ vs } k{+}1) > \thetaval \text{ with 95\% bootstrap CI lower bound} > \thetaval\}
\end{equation}

where $\thetaval = 0.60$ corresponds to Cohen's $h \approx 0.20$ (small effect size). We require the lower bound of a 95\% bootstrap confidence interval (10,000 resamples) to exceed $\thetaval$, ensuring statistical reliability.

We acknowledge that $\thetaval = 0.60$ is a parameter choice. Section~\ref{sec:theta} presents a full sensitivity analysis at $\thetaval \in \{0.55, 0.60, 0.65, 0.70\}$.

\subsection{Downstream Filtering Strategies}

We evaluate three filtering strategies for continued pretraining:
\begin{itemize}[nosep]
    \item \textbf{No filter}: use all data including contaminated text.
    \item \textbf{Binary}: remove all text classified as synthetic (any depth $> 0$).
    \item \textbf{Graduated}: weight text by estimated depth, down-weighting deeper generations with parameter $\alpha$.
\end{itemize}


%==============================================================================
\section{$\kstar$ Characterization Results}
\label{sec:results}
%==============================================================================

\subsection{C4 Web Text: $\kstar{=}2$ as Dominant Mode}
\label{sec:c4}

Table~\ref{tab:c4_main} presents $\kstar$ across 9 model configurations on C4 web text (5K documents, depth 0--5). Eight of nine configurations yield $\kstar{=}2$, with Mistral-7B as the sole exception at $\kstar{=}1$.

\begin{table}[t]
\centering
\caption{$\kstar$ determination on C4 web text across 9 model configurations. AUC values are 5-fold cross-validated means. CI lower denotes the 2.5th percentile of 10K bootstrap resamples. Bold indicates $\kstar{=}2$; $\dagger$ marks the rewrite mode pilot (not used for cross-model claims).}
\label{tab:c4_main}
\small
\begin{tabular}{llcccccc}
\toprule
Model & Mode & AUC 0v1 & AUC 1v2 & AUC 2v3 & CI$_{1v2}$ lower & $\kstar$ \\
\midrule
Qwen-1.5B-Inst$^\dagger$ & rewrite & 0.9956 & 0.6957 & 0.5893 & 0.6859 & \textbf{2} \\
Pythia-1.4B & cont & 0.9932 & 0.6539 & 0.5515 & 0.6454 & \textbf{2} \\
OLMo-1B & cont & 0.9970 & 0.6146 & 0.5497 & 0.6045 & \textbf{2} \\
Qwen2.5-7B & cont & 0.9644 & 0.6178 & 0.5460 & 0.6079 & \textbf{2} \\
LLaMA-3.1-8B & cont & 0.9225 & 0.6123 & 0.5274 & 0.6028 & \textbf{2} \\
Qwen-1.5B-Inst & cont & 0.9972 & 0.6114 & 0.5609 & 0.6015 & \textbf{2} \\
Qwen-1.5B base & cont & 0.9870 & 0.6262 & 0.5496 & 0.6168 & \textbf{2} \\
\midrule
Mistral-7B & cont & 0.9843 & 0.5935 & 0.5402 & 0.5835 & 1 \\
\bottomrule
\end{tabular}
\end{table}

Several observations merit attention. First, the 0v1 AUC is near-perfect ($>0.92$) for all models, confirming that first-generation synthetic text is readily distinguishable from human text. Second, the 1v2 AUC values cluster narrowly between 0.61 and 0.70 (excluding the rewrite mode pilot at 0.70), and all 2v3 AUC values fall below $\thetaval$. Third, confidence interval margins above $\thetaval$ are small for several models: Qwen-cont at +0.15pp, LLaMA-8B at +0.28pp, OLMo at +0.45pp. These borderline cases are a limitation we address in the $\thetaval$-sensitivity analysis (Section~\ref{sec:theta}).

The Mistral-7B exception ($\kstar{=}1$ on C4, with $P(\text{AUC} \leq 0.60) = 0.90$) is notable because this same model achieves $\kstar{=}3$ on arXiv (Section~\ref{sec:arxiv}), providing direct evidence that $\kstar$ depends on the source domain, not on model identity alone.


\subsection{Domain Hierarchy}
\label{sec:domains}

Table~\ref{tab:domain_summary} summarizes $\kstar$ across four source domains. The domain hierarchy progresses from Wikipedia ($\kstar{=}1$) through C4/Code ($\kstar \in \{1,2\}$) to arXiv ($\kstar{=}3$), though the intermediate domains show model-dependent rather than uniform behavior.

\begin{table}[t]
\centering
\caption{$\kstar$ summary across four source domains. Fractions indicate the proportion of tested model configurations achieving the stated $\kstar$ value. ``Scorer-controlled'' indicates cross-family scorer validation was applied.}
\label{tab:domain_summary}
\small
\begin{tabular}{lccl}
\toprule
Domain & $\kstar$ & Models & Note \\
\midrule
Wikipedia & 1 & 6/6 & Scorer-controlled \\
C4 web text & 2 (1) & 8/9 (1/9) & Mistral-7B is $\kstar{=}1$ \\
Code/Python & 2 (1) & 2/5 (3/5) & Model-dependent \\
arXiv & 3 & 5/5 & Domain-competent models only \\
\bottomrule
\end{tabular}
\end{table}

The domain hierarchy is cleanest at the extremes: Wikipedia uniformly yields $\kstar{=}1$ and arXiv uniformly yields $\kstar{=}3$ among domain-competent models. The intermediate domains (C4, Code) show a mixed pattern where some models reach $\kstar{=}2$ while others remain at $\kstar{=}1$, suggesting that domain structural complexity sets an upper bound that model-specific generation competence may or may not reach.


\subsection{arXiv: $\kstar{=}3$ Cross-Validated Across Five Models}
\label{sec:arxiv}

Table~\ref{tab:arxiv} presents the arXiv results. All five domain-competent models (spanning four architecture families) consistently yield $\kstar{=}3$, with the boundary pair (2v3) AUC ranging from 0.6215 to 0.6544.

\begin{table}[t]
\centering
\caption{arXiv domain $\kstar$ determination. Self-scorer and cross-family scorer results shown. OLMo-1B excluded (capacity failure; generated out-of-domain text).}
\label{tab:arxiv}
\small
\begin{tabular}{lcccccc}
\toprule
Model & AUC 0v1 & AUC 1v2 & AUC 2v3 & CI$_{2v3}$ lower & $\kstar$ & Cross-scorer $\kstar$ \\
\midrule
Qwen-1.5B & 0.9946 & 0.8580 & 0.6427 & 0.6275 & 3 & 3 \\
Pythia-1.4B & --- & --- & 0.6544 & 0.6401 & 3 & --- \\
Qwen-7B & --- & --- & 0.6255 & --- & 3 & 2\tablefootnote{Qwen-7B cross-scorer 2v3 AUC = 0.5989, below $\thetaval$.} \\
LLaMA-8B & --- & --- & 0.6215 & --- & 3 & 3 \\
Mistral-7B & --- & --- & 0.6275 & --- & 3 & 3 \\
\bottomrule
\end{tabular}
\end{table}

The cross-scorer validation reveals that Qwen-7B drops from $\kstar{=}3$ to $\kstar{=}2$ under cross-family scoring (2v3 AUC = 0.5989), yielding a 75\% preservation rate for arXiv cross-scorer results (3/4). This is lower than the 92\% (11/12) observed for C4, suggesting scorer independence is itself domain-dependent and never fully holds.

LaTeX markup contributes a measurable portion of the arXiv signal: removing LaTeX formatting reduces the 2v3 AUC from 0.6436 to 0.6229 ($\Delta = -2.1$pp), narrowing the margin above $\thetaval$ from +4.3pp to +2.3pp. While $\kstar{=}3$ is preserved after removal, more thorough structural normalization (citations, equation numbering, section headers) might further erode this margin.


\subsection{Code/Python Domain}
\label{sec:code}

Table~\ref{tab:code} presents results for the Code/Python domain. Unlike C4 (8/9 at $\kstar{=}2$) and arXiv (5/5 at $\kstar{=}3$), the Code domain shows a 2:3 split between $\kstar{=}2$ and $\kstar{=}1$.

\begin{table}[t]
\centering
\caption{Code/Python domain $\kstar$ determination. $\ddagger$ indicates preliminary worst-case estimates from running experiments. JSD ratio is JSD$_{1v2}$/JSD$_{2v3}$.}
\label{tab:code}
\small
\begin{tabular}{lccccc}
\toprule
Model & AUC 0v1 & AUC 1v2 & AUC 2v3 & JSD ratio & $\kstar$ \\
\midrule
Qwen-1.5B & 0.9388 & 0.7382 & 0.5300 & 11.7$\times$ & 2 \\
LLaMA-3.1-8B & 0.9161 & 0.7394 & 0.5112 & 13.4$\times$ & 2 \\
\midrule
Pythia-1.4B & 0.9592 & 0.5983 & 0.5215 & --- & 1 \\
Mistral-7B$^\ddagger$ & 0.91--0.95 & 0.52--0.58 & 0.50--0.53 & --- & 1 \\
Qwen-7B$^\ddagger$ & 0.93--0.96 & 0.54--0.60 & 0.50--0.54 & --- & 1 \\
\bottomrule
\end{tabular}
\end{table}

Qwen-1.5B and LLaMA-8B both show strong $\kstar{=}2$ with high 1v2 AUC ($>0.73$) and sharp JSD cliffs (11--13$\times$ drop at the 1v2$\to$2v3 boundary). These cliff dynamics differ from C4, where the sharpest drop occurs at 0v1$\to$1v2 rather than 1v2$\to$2v3, suggesting that code's syntactic regularity creates a distinct signal decay pattern even when $\kstar$ matches.

However, Pythia-1.4B achieves only $\kstar{=}1$ on Code (aggregate 1v2 AUC = 0.5983, with 3/5 folds above $\thetaval$ but mean below), and preliminary estimates for Mistral-7B and Qwen-7B also predict $\kstar{=}1$.
This 2/5 support for $\kstar{=}2$ means the Code domain does not support a clean ``ceiling$=$2'' interpretation.
Instead, code generation competence appears necessary to reach $\kstar{=}2$ in this domain, with models trained predominantly on non-code data (Pythia on The Pile) failing to produce sufficiently structured code to sustain depth discrimination.


\subsection{Scale Analysis}
\label{sec:scale}

Table~\ref{tab:scale} tracks the 0v1 AUC trend across Qwen model scales on C4, along with a preliminary estimate for 32B.

\begin{table}[t]
\centering
\caption{Scale trend for Qwen models on C4. $\ddagger$ indicates preliminary estimate. 0v1 AUC declining with scale reflects higher generation quality at larger scales.}
\label{tab:scale}
\small
\begin{tabular}{lcccc}
\toprule
Model & AUC 0v1 & AUC 1v2 & CI$_{1v2}$ lower & $\kstar$ \\
\midrule
Qwen-1.5B cont & 0.9972 & 0.6114 & 0.6015 & 2 \\
Qwen-7B & 0.9644 & 0.6178 & 0.6079 & 2 \\
Qwen-14B & 0.9598 & 0.6210 & 0.6119 & 2 \\
Qwen-32B$^\ddagger$ & 0.90--0.94 & 0.58--0.62 & 0.57--0.61 & 1--2 \\
\bottomrule
\end{tabular}
\end{table}

The 0v1 AUC decreases monotonically with scale (0.997 $\to$ 0.964 $\to$ 0.960), reflecting that larger models generate text closer to human quality. The 1v2 AUC shows no clear monotonic trend in the 1.5B--14B range, but the preliminary 32B estimate suggests a possible drop toward or below $\thetaval$. If confirmed, this would indicate that $\kstar{=}2$ for web text may be limited to models of 14B parameters or fewer, a significant constraint on the generalizability of our findings to frontier-scale models.


\subsection{Wikipedia: Scorer Confound and $\kstar{=}1$}
\label{sec:wiki}

Wikipedia presents a methodological cautionary tale. Under same-family scoring, Qwen generators yield $\kstar{=}2$ (AUC 1v2 = 0.6193). However, cross-family scorer validation (Table~\ref{tab:wiki_confound}) reveals this as a scorer artifact: all 6 models yield $\kstar{=}1$ under cross-family scoring.

\begin{table}[t]
\centering
\caption{Wikipedia scorer confound. Same-family scorer inflates $\kstar$ by +1 for Qwen generators. Cross-family scorer (Pythia) gives the corrected value.}
\label{tab:wiki_confound}
\small
\begin{tabular}{lccc}
\toprule
Condition & AUC 1v2 & $\Delta$ from same-scorer & $\kstar$ \\
\midrule
Qwen same-family scorer & 0.6193 & --- & 2 \\
Cross-family scorer (Pythia) & 0.5870 & $-$0.024 & 1 \\
\bottomrule
\end{tabular}
\end{table}

The Wikipedia scorer confound raises a general concern: if same-family scoring inflates $\kstar$ by +1 on Wikipedia, similar inflation might occur on other domains. We tested cross-scorer validation on C4 (11/12 preserved) and arXiv (3/4 preserved), finding smaller effects. However, the absence of detected confound on C4 does not guarantee its absence, as not all model$\times$domain combinations were tested.


\subsection{Cross-Model Transfer}
\label{sec:transfer}

Table~\ref{tab:transfer} presents the cross-model transfer analysis using a 4-generator $\times$ 4-scorer matrix.

\begin{table}[t]
\centering
\caption{Cross-model transfer summary. Retention is accuracy relative to within-model scoring. ``Cross avg PW-AUC'' is the mean pairwise AUC across all cross-model scoring conditions.}
\label{tab:transfer}
\small
\begin{tabular}{lcc}
\toprule
Transfer Condition & Retention & Note \\
\midrule
Pythia $\leftrightarrow$ OLMo & 91--92\% & Same pretraining paradigm \\
Qwen-cont $\to$ others & 80--85\% & Cross-family \\
Qwen-rewrite $\to$ others & 61--63\% & Mode-specific artifacts \\
\midrule
0v1 cross-model (all pairs) & $>0.96$ & Binary signal fully transfers \\
Within-model avg PW-AUC & 0.6099 & --- \\
Cross-model avg PW-AUC & 0.5905 & Below $\thetaval$ \\
\bottomrule
\end{tabular}
\end{table}

Binary detection (0v1) transfers perfectly across all model pairs ($>0.96$ AUC). However, cross-model average pairwise AUC for depth discrimination is 0.5905, below $\thetaval = 0.60$. This means that while binary contamination signals are fully model-independent, fine-grained depth discrimination does not reliably transfer across model families in aggregate.

\subsection{Cross-Scorer $\kstar$ Validation}
\label{sec:cross-scorer}

Table~\ref{tab:cross_scorer} presents the 4-generator $\times$ 4-scorer $\kstar$ matrix on C4.

\begin{table}[t]
\centering
\caption{Cross-scorer $\kstar$ matrix on C4 (4 generators $\times$ 4 scorers). Only one combination yields $\kstar{=}1$; the maximum scorer delta in AUC is 0.025.}
\label{tab:cross_scorer}
\small
\begin{tabular}{lcccc}
\toprule
& \multicolumn{4}{c}{Scorer} \\
\cmidrule(lr){2-5}
Generator & Pythia & OLMo & Qwen-base & Qwen-cont \\
\midrule
Pythia & $\kstar{=}2$ & $\kstar{=}2$ & $\kstar{=}2$ & --- \\
OLMo & $\kstar{=}2$ & $\kstar{=}2$ & $\kstar{=}2$ & --- \\
Qwen-base & $\kstar{=}2$ & $\kstar{=}2$ & $\kstar{=}2$ & --- \\
Qwen-cont & $\kstar{=}1$ & $\kstar{=}2$ & $\kstar{=}2$ & $\kstar{=}2$ \\
\bottomrule
\end{tabular}
\end{table}

11/12 generator$\times$scorer combinations yield $\kstar{=}2$, with the sole exception being Qwen-cont scored by Pythia (1v2 AUC = 0.5979, $\kstar{=}1$). Qwen-cont is also the most borderline case in Table~\ref{tab:c4_main} (CI lower = 0.6015, +0.15pp margin), and its failure under cross-scoring underscores the fragility of this particular determination.

This matrix does not include Mistral or LLaMA as generators, so the cross-scorer coverage is incomplete. Expanding to all model pairs would require approximately 30 additional experiments.


%==============================================================================
\section{Robustness and Ablations}
\label{sec:ablations}
%==============================================================================

\subsection{$\thetaval$-Sensitivity Analysis}
\label{sec:theta}

Table~\ref{tab:theta} shows how $\kstar$ determinations change across threshold values. The dependence on $\thetaval$ is the most significant robustness concern in our study.

\begin{table}[t]
\centering
\caption{C4 $\kstar{=}2$ count (out of 9 configurations) at different $\thetaval$ thresholds. ``Most marginal'' indicates the model with the smallest margin above $\thetaval$.}
\label{tab:theta}
\small
\begin{tabular}{lcl}
\toprule
$\thetaval$ & Count $\kstar{=}2$ & Most marginal \\
\midrule
0.55 & 9/9 & All well above \\
0.60 & 8/9 & Qwen-cont (+0.15pp) \\
0.65 & $\sim$5/9 & Multiple models flip \\
0.70 & $\sim$2/9 & Only Pythia, Qwen-rewrite survive \\
\bottomrule
\end{tabular}
\end{table}

At $\thetaval = 0.60$ (our default, corresponding to Cohen's $h \approx 0.20$ small effect size), 8/9 configurations yield $\kstar{=}2$.
At $\thetaval = 0.65$ (medium-small effect), approximately half flip to $\kstar{=}1$.
At $\thetaval = 0.70$, only 2/9 retain $\kstar{=}2$.
Permutation tests confirm all 1v2 AUC values are significantly above 0.50 ($p < 0.001$), establishing that the depth signal is real (non-random). The practical question is whether the effect size is ``meaningful,'' which is inherently context-dependent.


\subsection{Decoding Strategy Ablation}
\label{sec:decoding}

Table~\ref{tab:decoding} shows that $\kstar{=}2$ is invariant to decoding strategy for both Qwen-1.5B and Pythia-1.4B on C4.

\begin{table}[t]
\centering
\caption{Decoding strategy ablation on C4 (Qwen-1.5B rewrite and Pythia-1.4B continuation). All configurations yield $\kstar{=}2$.}
\label{tab:decoding}
\small
\begin{tabular}{lcccc}
\toprule
& \multicolumn{2}{c}{Qwen-1.5B (rewrite)} & \multicolumn{2}{c}{Pythia-1.4B (cont)} \\
\cmidrule(lr){2-3} \cmidrule(lr){4-5}
Decoding & AUC 1v2 & $\kstar$ & AUC 1v2 & $\kstar$ \\
\midrule
Greedy & 0.6892 & 2 & 0.8838 & 2 \\
Nucleus $p{=}0.9$ & 0.6862 & 2 & 0.6770 & 2 \\
Temp=0.7 & 0.6885 & 2 & 0.7265 & 2 \\
Temp=1.2 & 0.7037 & 2 & 0.6172 & 2 \\
\bottomrule
\end{tabular}
\end{table}

The Qwen results show tight AUC clustering (range 0.69--0.70). Pythia shows much larger variation, with greedy decoding at 0.8838 (anomalously high due to repetition collapse in Pythia's greedy output) and temperature 1.2 at 0.6172 (near the borderline). The Pythia greedy anomaly suggests that part of the measured ``depth signal'' in that condition reflects text degeneracy artifacts rather than genuine depth information.


\subsection{Text Length Ablation}
\label{sec:length}

Table~\ref{tab:length} shows the effect of text segment length on $\kstar$ determination.

\begin{table}[t]
\centering
\caption{Text length ablation (Qwen continuation, C4). The 512-token condition fails bootstrap CI at $\thetaval{=}0.60$.}
\label{tab:length}
\small
\begin{tabular}{lcccc}
\toprule
Length (tokens) & AUC 1v2 & CI$_{1v2}$ lower & Verdict & $\kstar$ \\
\midrule
128 & 0.6219 & --- & Pass & 2 \\
256 & 0.6282 & --- & Pass & 2 \\
512 & 0.6073 & 0.5979 & \textsc{borderline} & 2\tablefootnote{Point estimate passes but CI lower = 0.5979 $< \thetaval$.} \\
\bottomrule
\end{tabular}
\end{table}

The 512-token condition is borderline: the point estimate (0.6073) exceeds $\thetaval$ but the CI lower bound (0.5979) does not. This is concerning for practical applicability, since real documents are typically longer than 256 tokens. The finding suggests $\kstar{=}2$ may be most reliable for short text segments, limiting ecological validity.


\subsection{Base vs.\ Instruct Models}

Qwen-1.5B base (AUC 1v2 = 0.6262, CI lower = 0.6168) yields the same $\kstar{=}2$ as Qwen-1.5B-Instruct continuation (0.6114, CI lower = 0.6015), with overlapping confidence intervals. Instruction tuning (alignment) does not significantly affect $\kstar$ on C4.


\subsection{Feature Analysis}
\label{sec:features}

Cross-model feature ranking shows Spearman $\rho > 0.85$, with surprisal-based features (surp\_d1\_std, surp\_std, surp\_d2\_std) consistently ranked as top discriminators across all tested models. However, the 1v2 depth pair, which is critical for $\kstar{=}2$ determination, shows the weakest cross-model feature consistency ($\tau \sim 0.37$), meaning that the exact features driving the boundary-pair classification vary across models.

\subsection{Neural Probe Analysis}
\label{sec:neural}

Table~\ref{tab:neural} compares $\kstar$ obtained from frozen neural representations (self-embeddings + MLP probe) against the 15D model-free features.

\begin{table}[t]
\centering
\caption{Neural probe vs.\ 15D model-free features on C4. Neural $\kstar$ is model-dependent; 15D features give consistent results.}
\label{tab:neural}
\small
\begin{tabular}{lccc}
\toprule
Model & $\kstar$ (15D) & $\kstar$ (neural) & Note \\
\midrule
Pythia-1.4B & 2 & 2 & Agreement \\
OLMo-1B & 2 & 2 & Agreement \\
Qwen-1.5B-Inst & 2 & 3 & Neural $>$ 15D \\
\bottomrule
\end{tabular}
\end{table}

Qwen's neural probe achieves $\kstar{=}3$ (2v3 AUC = 0.6159, CI lower = 0.6099), indicating that richer representations can extract additional depth information beyond what 15D features capture. This has two interpretations: (1) the model-free features are practical because they give consistent $\kstar$ across models, or (2) $\kstar{=}2$ from 15D features is a measurement floor rather than a data-intrinsic ceiling. We adopt interpretation (1) for our headline results but acknowledge (2) as a valid concern.


\subsection{Retrain Chain}
\label{sec:retrain}

Table~\ref{tab:retrain} contrasts prompt-chain and retrain-chain $\kstar$ values.

\begin{table}[t]
\centering
\caption{Prompt-chain vs.\ retrain-chain (LoRA) $\kstar$ comparison. Retrain chains exhibit much higher $\kstar$ due to accumulated model adaptation artifacts.}
\label{tab:retrain}
\small
\begin{tabular}{llcccc}
\toprule
Model & Chain & AUC 1v2 & AUC 3v4 & AUC 6v7 & $\kstar$ \\
\midrule
Pythia-1.4B & Prompt & 0.6539 & --- & --- & 2 \\
Pythia-1.4B & Retrain & 0.9484 & 0.7982 & 0.6340 & $\geq$4 \\
Qwen-1.5B & Prompt & 0.6114 & --- & --- & 2 \\
Qwen-1.5B & Retrain & --- & --- & 0.6410 & 7 \\
\bottomrule
\end{tabular}
\end{table}

Retrain chains extend $\kstar$ to 7 (both Qwen 6v7 = 0.6410, CI lower = 0.6376, and Pythia 6v7 = 0.6340, CI lower = 0.6275), far exceeding the prompt-chain $\kstar{=}2$. This large gap (3.5$\times$) demonstrates that $\kstar$ depends heavily on the contamination mechanism. Retrain chains accumulate LoRA adaptation artifacts (vocabulary narrowing, repetition bias) that create persistent detectable signatures across many generations. The prompt-chain $\kstar{=}2$ reflects text-intrinsic information loss only, making it the more conservative and generalizable estimate.


\subsection{Subsampling Power Analysis}

The minimum sample size for 100\% detection of $\kstar{=}2$ is $n = 2{,}000$ documents. At $n = 1{,}000$, only 30\% of bootstrap resamples yield significant $\kstar{=}2$ determinations. At $n = 5{,}000$, the mean 1v2 AUC is 0.6288. The $n = 2{,}000$ requirement reflects the small effect size: depth discrimination is a statistical property of document collections, not a per-document signal.


%==============================================================================
\section{Depth Estimation}
\label{sec:estimation}
%==============================================================================

We evaluate ordinal depth estimation as a practical application of $\kstar$ characterization.

\subsection{Ordinal vs.\ Multiclass vs.\ LLM Baselines}

Table~\ref{tab:estimation} compares four approaches to 5-class (depth 0--4) or 6-class (depth 0--5) classification.

\begin{table}[t]
\centering
\caption{Depth estimation accuracy comparison. OBD = ordinal binary decomposition with LightGBM. Random baseline for 6-class is 16.7\%.}
\label{tab:estimation}
\small
\begin{tabular}{lcc}
\toprule
Method & Accuracy & Note \\
\midrule
OBD-LightGBM (15D) & 42.3\% & 6-class, best ordinal method \\
Multiclass LightGBM & 41.8\% & +0.47pp for OBD \\
Ordinal logistic regression & 36.4\% & Linear baseline \\
GPT-5.5 zero-shot & 33.2\% & 5-class, 500 samples \\
Claude Opus 4.6 zero-shot & 33.1\% & 5-class, 500 samples \\
OBD (on same 500 samples) & 38.6\% & Matched evaluation \\
\midrule
Binary detection (0v1) & $\sim$88\% & All methods \\
\bottomrule
\end{tabular}
\end{table}

The OBD advantage over multiclass is minimal (+0.47pp), and the difference between OBD and frontier LLMs is not statistically significant (McNemar $p = 0.081$ vs.\ GPT-5.5, $p = 0.069$ vs.\ Claude). All fine-grained depth estimation methods achieve modest accuracy ($< 43\%$), while binary detection (depth 0 vs.\ $>0$) achieves approximately 88\% accuracy across all methods. This gap illustrates the fundamental difficulty of fine-grained depth estimation relative to binary detection.


\subsection{Binary Detector Comparison}

Existing binary detectors (DivEye, Binoculars) used as depth proxies via confidence score correlation achieve $\kstar{=}1$ (unable to discriminate beyond first generation), maximum Spearman $\rho = 0.46$ with true depth, and 6-class accuracy $\leq 19.9\%$. The OBD approach achieves $\kstar{=}2$, demonstrating that the 15D features capture depth information beyond what binary detector scores provide. However, this comparison involves a 15-feature trained classifier against single-score methods not designed for depth estimation.

\subsection{Ordinal Accuracy}

Pythia-1.4B achieves ordinal accuracy of 38.3\% on 6-class depth estimation (random baseline 16.7\%), and OLMo-1B achieves 36.9\%. These values, while above random, indicate that fine-grained depth assignment remains challenging at the individual document level.


%==============================================================================
\section{Downstream Filtering}
\label{sec:filtering}
%==============================================================================

\subsection{C4 Filtering}

Table~\ref{tab:filtering_c4} presents downstream evaluation of filtering strategies on continued pretraining with Qwen-1.5B (LoRA, 5 seeds).

\begin{table}[t]
\centering
\caption{C4 downstream filtering results (Qwen-1.5B, 30\% contamination, 5 seeds). Base model MMLU = 0.6018 (no continued pretraining). All continued pretraining degrades below base.}
\label{tab:filtering_c4}
\small
\begin{tabular}{lcc}
\toprule
Strategy & MMLU & HellaSwag \\
\midrule
Base model (no CPT) & 0.6018 & --- \\
\midrule
No filter & 0.5865 $\pm$ 0.0010 & --- \\
Binary-ours & \textbf{0.5941 $\pm$ 0.0008} & --- \\
Graduated ($\alpha{=}0.9$) & 0.5892 $\pm$ 0.0008 & --- \\
DSIR & $\sim$0.5867 & --- \\
Surprisal-based & 0.5894 & --- \\
\bottomrule
\end{tabular}
\end{table}

Binary filtering achieves the best MMLU (0.5941, $\Delta = +0.76$pp over no filter) and HellaSwag ($\Delta = +0.63$pp). Graduated filtering performs between binary and no-filter but provides no significant improvement over simpler baselines. All continued pretraining conditions degrade below the base model (0.6018 MMLU), meaning the best filtering strategy still results in net degradation from contaminated data.

\subsection{Dose-Response Analysis}

Table~\ref{tab:dose} presents filtering performance across three contamination ratios.

\begin{table}[t]
\centering
\caption{Dose-response: binary filtering advantage (in HellaSwag pp) over no-filter at three contamination ratios. Graduated filtering equals no-filter at all ratios.}
\label{tab:dose}
\small
\begin{tabular}{lccc}
\toprule
Ratio & Binary $\Delta$ (HellaSwag) & Graduated $\Delta$ & $p$-value (binary) \\
\midrule
10\% & +0.38pp & $\approx$0 & $\leq$0.004 \\
30\% & +0.61pp & $\approx$0 & $\leq$0.004 \\
50\% & +0.51pp & $\approx$0 & $\leq$0.004 \\
\bottomrule
\end{tabular}
\end{table}

Binary filtering maintains a statistically significant advantage at all contamination ratios ($p \leq 0.004$ on HellaSwag). Graduated filtering collapses to no-filter performance at \emph{all} ratios. If depth-aware filtering were providing genuine value, we would expect its advantage to increase with contamination severity. The absence of any graduated benefit at any ratio is a strong negative result for depth-aware curation.

\subsection{Pythia Replication}

Table~\ref{tab:pythia_filter} shows that the C4 filtering result does not replicate on Pythia-1.4B.

\begin{table}[t]
\centering
\caption{Pythia-1.4B filtering replication (LoRA rank=8, 3 seeds). All strategies produce indistinguishable results.}
\label{tab:pythia_filter}
\small
\begin{tabular}{lcc}
\toprule
Strategy & MMLU & Conclusion \\
\midrule
Binary-ours & 0.2585 $\pm$ 0.0004 & --- \\
No filter & 0.2568 $\pm$ 0.0012 & --- \\
$\Delta_{\max}$ & 0.0017 & Within noise \\
\bottomrule
\end{tabular}
\end{table}

All three strategies (binary, graduated, no filter) produce results within noise ($\Delta_{\max} = 0.17$pp), indicating that the filtering benefit observed for Qwen does not generalize to Pythia. The Qwen-specific result may reflect a particular sensitivity of Qwen-1.5B's tokenizer or architecture to binary filtering, rather than a general principle.

\subsection{arXiv Filtering}

Table~\ref{tab:arxiv_filter} presents the arXiv domain filtering results, where $\kstar{=}3$ would predict the strongest case for depth-aware curation.

\begin{table}[t]
\centering
\caption{arXiv filtering (Qwen-7B, 5 seeds). Despite $\kstar{=}3$, no strategy differs from no-filter.}
\label{tab:arxiv_filter}
\small
\begin{tabular}{lcc}
\toprule
Strategy & MMLU & Conclusion \\
\midrule
Graduated & 0.5900 $\pm$ 0.0011 & --- \\
Binary & 0.5876 $\pm$ 0.0004 & --- \\
No filter & 0.5899 $\pm$ 0.0008 & --- \\
\bottomrule
\end{tabular}
\end{table}

All strategies produce indistinguishable results on arXiv. This is the most damaging negative result for depth-aware curation: if $\kstar{=}3$ represents the domain with the richest depth information, and even here graduated filtering provides no benefit, then depth-aware curation has no demonstrated practical advantage on any tested domain.


\subsection{Statistical Tests}
\label{sec:stats}

Holm-Bonferroni corrected paired tests across all filtering comparisons yield 8/15 significant pairs on MMLU and 9/15 on HellaSwag. Binary-ours vs.\ all alternatives is significant on MMLU. Graduated ($\alpha{=}0.9$) vs.\ DSIR yields corrected $p = 0.0523$ (not significant), meaning graduated filtering cannot be distinguished from DSIR, a method that uses no depth information.

High Cohen's $d$ values ($>8$) reflect extremely low within-seed variance rather than practically large effects. The absolute deltas remain below 1pp MMLU.


%==============================================================================
\section{Discussion}
\label{sec:discussion}
%==============================================================================

\subsection{The Rapid Convergence Finding}

The central result of this study is that generational discrimination saturates within three passes ($\kstar \leq 3$) across all tested conditions. The DPI guarantees that $\kstar$ must be finite, but provides no upper bound. Finding $\kstar \leq 3$ empirically means that recursive generation erases distinguishable structure extremely quickly. For web text, the ``one-wash laundering'' effect (where a single additional pass beyond $\kstar{=}2$ renders text indistinguishable from deeper generations) has direct implications for data governance: there is limited value in estimating exact depth beyond the first 1--2 generations.

\subsection{Domain Dependence: Clear at Extremes, Messy in the Middle}

The domain hierarchy is robust at its endpoints (Wikipedia uniformly $\kstar{=}1$, arXiv uniformly $\kstar{=}3$) but model-dependent in intermediate-complexity domains. C4 shows 8/9 at $\kstar{=}2$ with one exception, while Code shows only 2/5 at $\kstar{=}2$. This mixed evidence in intermediate domains means the ``domain ceiling'' framework is descriptive rather than predictive: domain structure may set an upper bound, but whether individual models reach that bound depends on their domain-specific generation competence.

\subsection{The C3 Negative Result}

Our downstream filtering experiments yield a clear negative result for depth-aware curation. Binary filtering provides a small, model-specific benefit (Qwen-1.5B on C4 only), but graduated depth-aware weighting provides \emph{zero} benefit at any contamination level, on any domain, for any model tested. Combined with the Pythia replication failure and the arXiv negative result, the practical utility of $\kstar$ for data curation remains undemonstrated. We report this honestly as a negative finding rather than suppressing it.

\subsection{Limitations}

\paragraph{Threshold dependence.} Our $\kstar$ determinations depend sensitively on $\thetaval$. At $\thetaval{=}0.60$ (small effect), 8/9 C4 configurations yield $\kstar{=}2$; at $\thetaval{=}0.65$, approximately half flip to $\kstar{=}1$. Whether $\thetaval{=}0.60$ represents ``meaningful'' discrimination is context-dependent.

\paragraph{Scale generalization.} Our experiments cover 1B--14B parameter models. The declining 0v1 AUC trend with scale and the borderline 32B preliminary estimate suggest $\kstar$ may decrease at frontier scales ($>$32B), limiting the relevance of our findings to the most capable current models.

\paragraph{Scorer confound.} Wikipedia revealed that same-family scoring can inflate $\kstar$ by +1. While cross-scorer validation on C4 and arXiv showed smaller effects, not all model$\times$domain combinations were tested. The incomplete coverage leaves open the possibility of undetected confounds.

\paragraph{Contamination mechanism.} We study only prompt-chain generation. The 3.5$\times$ difference between prompt-chain and retrain-chain $\kstar$ (2 vs.\ 7) indicates mechanism dependence. Real-world contamination involves diverse mechanisms (paraphrasing, translation, summarization, multi-model pipelines) not captured in our controlled setup.

\paragraph{Naturalistic validation.} We have not validated our laboratory $\kstar$ measurements against real-world contaminated corpora (e.g., Common Crawl temporal snapshots). Without ground-truth depth labels in the wild, ecological validity remains an open question.


%==============================================================================
\section{Conclusion}
%==============================================================================

We introduced generational discrimination depth $\kstar$ and provided the first systematic empirical characterization of this quantity across diverse models, domains, and conditions.
The core finding, that $\kstar$ saturates at $\leq 3$ across all tested conditions, establishes that recursive generation erases distinguishable depth structure rapidly.
Domain structural complexity emerges as the primary factor: simple domains (Wikipedia) saturate at $\kstar{=}1$, general web text (C4) at $\kstar{=}2$, and structured academic text (arXiv) at $\kstar{=}3$, though intermediate domains show model-dependent behavior.

Our downstream experiments provide an important negative result: depth-aware curation does not improve upon binary filtering in any tested condition.
This suggests that while $\kstar$ characterization advances our understanding of recursive contamination, translating this understanding into practical curation benefits requires further research, particularly at larger model scales and with more diverse contamination mechanisms.


%==============================================================================
% References
%==============================================================================
\bibliographystyle{plainnat}

\begin{thebibliography}{20}

\bibitem[Shumailov et~al.(2024)]{shumailov2024ai}
Ilia Shumailov, Zakhar Shumaylov, Yiren Zhao, Nicolas Papernot, Ross Anderson, and Yarin Gal.
\newblock AI models collapse when trained on recursively generated data.
\newblock \emph{Nature}, 631:755--759, 2024.

\bibitem[Alemohammad et~al.(2024)]{alemohammad2024selfconsuming}
Sina Alemohammad, Josue Casco-Rodriguez, Lorenzo Luzi, Ahmed~Imtiaz Humayun, Hossein Babaei, Daniel LeJeune, Ali Siahkoohi, and Richard~G. Baraniuk.
\newblock Self-consuming generative models go {MAD}.
\newblock In \emph{ICLR}, 2024.

\bibitem[Dohmatob et~al.(2024)]{dohmatob2024tale}
Elvis Dohmatob, Yunzhen Feng, and Julia Kempe.
\newblock A tale of tails: Model collapse as a change of scaling laws.
\newblock In \emph{ICML}, 2024.

\bibitem[Mitchell et~al.(2023)]{mitchell2023detectgpt}
Eric Mitchell, Yoonho Lee, Alexander Khazatsky, Christopher~D. Manning, and Chelsea Finn.
\newblock {DetectGPT}: Zero-shot machine-generated text detection using random perturbations.
\newblock In \emph{ICML}, 2023.

\bibitem[Drayson et~al.(2025)]{drayson2025preventing}
Benjamin Drayson, Thom Lake, and Yarin Gal.
\newblock Preventing model collapse in data flywheel with binary detection resampling.
\newblock \emph{arXiv preprint arXiv:2503.XXXXX}, 2025.

\bibitem[Xie et~al.(2023)]{xie2023data}
Sang~Michael Xie, Shibani Santurkar, Tengyu Ma, and Percy Liang.
\newblock Data selection for language models via importance resampling.
\newblock In \emph{NeurIPS}, 2023.

\bibitem[Cover and Thomas(2006)]{cover2006elements}
Thomas~M. Cover and Joy~A. Thomas.
\newblock \emph{Elements of Information Theory}.
\newblock Wiley-Interscience, 2nd edition, 2006.

\end{thebibliography}


%==============================================================================
\clearpage
\appendix
\section{Supplementary Experiments}
%==============================================================================

\subsection{Rewrite Mode Pilot (MVP)}
\label{app:mvp}

The initial feasibility study used Qwen2.5-1.5B-Instruct in rewrite mode on C4 (5K documents, depth 0--5). Results: $\kstar{=}2$ with AUC 1v2 = 0.6957, CI [0.6859, 0.7041], and JSD$_{0v1}$ = 0.2669, JSD$_{1v2}$ = 0.0163. While the highest AUC 1v2 among all C4 conditions, cross-model transfer retention of 61--63\% (vs.\ 80--92\% for continuation mode) indicates that the elevated signal reflects instruction-following style artifacts rather than transferable depth information.


\subsection{Pythia Decoding Anomaly}
\label{app:pythia_decoding}

Pythia-1.4B under greedy decoding produces AUC 1v2 = 0.8838, approximately 0.20 above the next-highest decoding condition (T=0.7 at 0.7265). Greedy decoding with Pythia triggers repetition collapse (the model generates highly repetitive sequences), creating artificial depth signatures that inflate discriminability. Temperature 1.2 yields the lowest AUC (0.6172), approaching the borderline. This 0.27 AUC range across decoding strategies (vs.\ 0.02 for Qwen) reflects Pythia's sensitivity to text degeneracy modes.


\subsection{Ordinal Accuracy Details}
\label{app:ordinal}

Pythia ordinal accuracy = 38.3\% (6-class, random = 16.7\%), OLMo = 36.9\%. While above random, these values indicate that fine-grained depth assignment at the individual document level remains unreliable.


\subsection{Subsampling Power Curve}
\label{app:power}

\begin{table}[h]
\centering
\caption{C4 subsampling power analysis ($\times$10 resamples per $n$).}
\label{tab:power}
\small
\begin{tabular}{lcc}
\toprule
$n$ (documents) & Detection rate (\%) & Mean AUC 1v2 \\
\midrule
1,000 & 30 & --- \\
2,000 & 100 & --- \\
5,000 & 100 & 0.6288 \\
\bottomrule
\end{tabular}
\end{table}


\subsection{LaTeX Markup Removal on arXiv}
\label{app:latex}

\begin{table}[h]
\centering
\caption{arXiv $\kstar$ with and without LaTeX markup.}
\label{tab:latex_removal}
\small
\begin{tabular}{lccc}
\toprule
Condition & AUC 2v3 & $\Delta$ & $\kstar$ \\
\midrule
Original (with LaTeX) & 0.6436 & --- & 3 \\
LaTeX stripped & 0.6229 & $-$0.021 & 3 \\
\bottomrule
\end{tabular}
\end{table}

$\kstar{=}3$ is preserved after LaTeX removal, but the margin above $\thetaval$ narrows from +4.4pp to +2.3pp.


\subsection{Qwen-1.5B Retrain Chain Extended to Depth 7}
\label{app:retrain_d7}

Qwen-1.5B retrain chain (LoRA) reaches $\kstar{=}7$ with 6v7 AUC = 0.6410, CI lower = 0.6376. Pythia-1.4B retrain chain yields consistent $\kstar{=}7$ (6v7 AUC = 0.6340, CI lower = 0.6275). The cross-family agreement ($\kstar{=}7$ for both) occurs despite approximately 3$\times$ difference in LoRA effective capacity, suggesting the retrain-chain ceiling is determined by the number of LoRA adaptation rounds rather than model capacity.


\subsection{arXiv Cross-Scorer Validation Details}
\label{app:arxiv_cross}

\begin{table}[h]
\centering
\caption{arXiv self-scorer vs.\ cross-family scorer $\kstar$.}
\label{tab:arxiv_cross}
\small
\begin{tabular}{lcc}
\toprule
Model & Self-scorer $\kstar$ & Cross-scorer $\kstar$ \\
\midrule
Qwen-1.5B & 3 & 3 \\
Qwen-7B & 3 & 2 (2v3 = 0.5989) \\
LLaMA-8B & 3 & 3 \\
Mistral-7B & 3 & 3 \\
\bottomrule
\end{tabular}
\end{table}

Cross-scorer preservation rate: 3/4 = 75\%, lower than C4 (11/12 = 92\%).


\subsection{Holm-Bonferroni Corrected Filtering Tests}
\label{app:stats}

\begin{table}[h]
\centering
\caption{Summary of Holm-Bonferroni corrected paired tests for C4 filtering.}
\label{tab:holm}
\small
\begin{tabular}{lcc}
\toprule
Metric & Significant pairs & Note \\
\midrule
MMLU & 8/15 & Binary vs.\ all significant \\
HellaSwag & 9/15 & Binary vs.\ all significant \\
Graduated vs.\ DSIR (MMLU) & --- & Corrected $p = 0.0523$ (n.s.) \\
\bottomrule
\end{tabular}
\end{table}


\subsection{OBD vs.\ LLM McNemar Tests}
\label{app:mcnemar}

On the matched 500-sample evaluation: OBD accuracy = 38.0\%, GPT-5.5 = 33.2\%, Claude Opus 4.6 = 33.1\%. McNemar tests: OBD vs.\ GPT-5.5 $p = 0.081$, OBD vs.\ Claude $p = 0.069$. Neither is significant at $\alpha = 0.05$, meaning the specialized 15-feature classifier does not significantly outperform zero-shot frontier LLMs on depth estimation.


\subsection{Aggregate Summary}
\label{app:aggregate}

\begin{table}[h]
\centering
\caption{Aggregate worst-case summary across all experiments.}
\label{tab:aggregate}
\small
\begin{tabular}{ll}
\toprule
Metric & Value \\
\midrule
Domain hierarchy clean fit & 78--85\% \\
Code ceiling$=$2 support & 2/5 models \\
C4 $\kstar{=}2$ at $\thetaval{=}0.60$ & 8/9 \\
C4 $\kstar{=}2$ at $\thetaval{=}0.65$ & $\sim$5/9 \\
Cross-scorer preservation rate & 65--75\% \\
C3 validated scope & Qwen-1.5B C4 binary only \\
C3 depth-aware advantage & Zero (graduated $=$ no-filter) \\
\bottomrule
\end{tabular}
\end{table}


\end{document}
